\documentclass[11pt,a4paper]{article}
\usepackage[margin=2cm]{geometry}
\usepackage[T1]{fontenc}
\usepackage{lmodern}
\usepackage{amsmath,amssymb,amsthm}
\usepackage{graphicx,booktabs,xcolor,listings,microtype}
\usepackage[font=small,labelfont=bf,skip=3pt]{caption}
% Compact spacing (the report has a 4-page limit)
\setlength{\parindent}{0pt}
\setlength{\parskip}{2.4pt plus 1pt}
\setlength{\abovedisplayskip}{4pt plus 1pt minus 1pt}
\setlength{\belowdisplayskip}{4pt plus 1pt minus 1pt}
\setlength{\abovedisplayshortskip}{2pt}
\setlength{\belowdisplayshortskip}{3pt}
\setlength{\textfloatsep}{6pt plus 2pt}
\setlength{\floatsep}{5pt plus 2pt}
\setlength{\intextsep}{5pt plus 2pt}
% Theorems, lemmas, facts and definitions share one counter
\newtheoremstyle{tight}{2.4pt}{0pt}{\itshape}{}{\bfseries}{.}{.5em}{}
\newtheoremstyle{tightrm}{2.4pt}{0pt}{}{}{\bfseries}{.}{.5em}{}
\theoremstyle{tight}
\newtheorem{thm}{Theorem}
\newtheorem{lem}[thm]{Lemma}
\newtheorem{fact}[thm]{Fact}
\theoremstyle{tightrm}
\newtheorem{defn}[thm]{Definition}
% Run-in section headings, numbered from 0
\newcounter{sect}\setcounter{sect}{-1}
\newcounter{subsect}[sect]
\renewcommand{\thesubsect}{\thesect.\arabic{subsect}}
\newcommand{\sect}[1]{\par\vspace{1.5pt}\noindent\stepcounter{sect}\textbf{\thesect. #1}\enspace}
\newcommand{\subsect}[1]{\par\vspace{1.5pt}\noindent\stepcounter{subsect}\textbf{\thesubsect\ #1}\enspace}
% Numbered methods (M1), (M2), ... in the numerical section
\newcounter{meth}\renewcommand{\themeth}{M\arabic{meth}}
\newcommand{\meth}{\stepcounter{meth}(\themeth)}
\newcommand{\R}{\mathbb{R}}
\newcommand{\Prob}{\mathbb{P}}
\newcommand{\e}{\mathrm{e}}
\newcommand{\one}{\mathbf{1}}
% Proofs without extra vertical space
\renewenvironment{proof}[1][Proof]{\par\textit{#1.}\enspace}{\qed\par}
% References printed as one compact paragraph
\renewenvironment{thebibliography}[1]
  {\par\vspace{1.5pt}\noindent\textbf{\refname}\enspace\footnotesize\setcounter{enumiv}{0}%
   \def\bibitem##1{\ifnum\value{enumiv}>0 \unskip\hskip 0.7em plus 0.2em minus 0.2em\fi
     \stepcounter{enumiv}[\arabic{enumiv}]~\ignorespaces}}
  {\par}
\lstset{basicstyle=\ttfamily\fontsize{6.6}{7.7}\selectfont,breaklines=true,columns=fullflexible, keepspaces=true,showstringspaces=false,upquote=true}
\begin{document}
\begin{center}
{\Large\bfseries How fast does the binomial model converge to Black-Scholes?}\\[1pt]
{\large\bfseries The order of convergence for calls and for discontious payoffs}
\end{center}
\vspace{-7pt}
\sect{Background and Plan.}
The binomial model of Cox, Ross and Rubinstein (CRR) values a European option with a finite sum. For increasing $n$ number of time periods, the option price of the binomial model converges to the black-Scholes model and it was obtained in [1, Sect. 7.4] by the central limit theorem (CLT). The CLT proves the convergence theorem, but it has no rate, which is necessary in practice, as n is the cost of computation. The convergence is also not monotone since binomial prices oscillate around the limit. If n is doubled, the rate can make the error larger (Section 4).
I prove that the the error is $O(1/n)$ for a call option, using the order $n^{-1/2}$ for the payoff with jump at the strike while including a oscillating leading term. The result is based of the expansion of the probability in binomial distributions. Next, I test from first principles the predictions using the cited formula for the coefficient of $1/n$. Section 5 tests the continutity assumption and section 6 justifies the assuptions I make.
\emph{Key idea:} The core of my analysis is that every price combines the binomial tail probabilities and that a normal approximation with continuity correction shows that each tail carries the size $n^{-1/2}$. This records where the strike sits between two nodes of thr binomial tree, and for a continuous payoff, the two tails cancel each other out, leaving a $n^{-1}$ error. For a discontinuous payoff these terms cancel exactly.

\sect{Setting and definitions.} 
\emph{Requirements:} I assume the probability topics up to the CLT, Taylor's theorem, and from Financial Mathematics~I [1] the multi-period binomial model, its risk-neutral pricing formula and the Black-Scholes formula. I define $\varphi(t)=\e^{-t^2/2}/\sqrt{2\pi}$ and $\Phi(t)=\int_{-\infty}^{t}\varphi$, and $x_n=O(y_n)$ means $|x_n|\le My_n$ for all $n\ge n_0$, where
$M$ and $n_0$ depend only on the fixed parameters $S,K,\sigma,T>0$, $r\in\R$ (initial share price, strike, volatility, expiry, continuously compounded interest rate respectively) and on $c_1,c_2$ below.

\begin{defn}[CRR model]
For $n> r^2T/\sigma^2$ put $h=T/n$, $u=\e^{\sigma\sqrt h}$, $d=u^{-1}$. In each of $n$ periods of length $h$, the share price is multiplied by $u$ or by $d$. Then $d<\e^{rh}<u$ (no arbitrage), and we set
\[
p_n=\frac{\e^{rh}-d}{u-d},\qquad q_n=1-p_n,\qquad \hat p_n=p_nu\,\e^{-rh}.
\]

Under the risk-neutral distribution the price at expiry is $Su^{J}d^{\,n-J}$ with $J\sim\mathrm{Bin}(n,p_n)$ [Prop.~6.24][1]. Since $p_nu+q_nd=\e^{rh}$, we have $1-\hat p_n=q_nd\,\e^{-rh}$, so $\hat p_n\in(0,1)$.
\end{defn}

\begin{defn}[Payoffs and prices]
For $c_1,c_2\in\R$ let $g(s)=(c_1s+c_2)\one_{\{s>K\}}$; its \emph{jump} at the strike is $[g]=c_1K+c_2$.
The $n$-step binomial price of the European option paying $g(S(T))$ at time $T$ is $V_n(g)=\e^{-rT}\,\mathbb{E}\,g(Su^Jd^{\,n-J})$ [Thm~6.25][1], and its Black--Scholes price is
\[
V(g)=c_1S\,\Phi(d_1)+c_2\,\e^{-rT}\Phi(d_2),\qquad
d_{1,2}=\frac{\ln(S/K)+(r\pm\sigma^2/2)T}{\sigma\sqrt T}
\]
[proof of Thm~7.3, eqs (26)-(27)][1]. The \emph{call} is $c_1=1,c_2=-K$, with payoff $(s-K)^+$, $[g]=0$
and prices $C_n$, $C$; the \emph{digital call} is $c_1=0,c_2=1$, with payoff $\one_{\{s>K\}}$, $[g]=1$ and prices $D_n$, $D$.
\end{defn}

\begin{defn}[Position of the strike in the lattice]
Let $\alpha_n=\tfrac n2+\dfrac{\ln(K/S)}{2\sigma\sqrt h}$. Since $Su^jd^{\,n-j}=S\e^{(2j-n)\sigma\sqrt h}$, the
node $j$ is in the money, $Su^jd^{\,n-j}>K$, exactly when $j>\alpha_n$. Let $a_n=\lfloor\alpha_n\rfloor+1$ be
the first such node, and put $\theta_n=a_n-\alpha_n\in(0,1]$ and $\Delta_n=1-2\theta_n\in[-1,1)$.
\end{defn}
So $\Delta_n=0$ when $K$ is midway (in log-price) between two neighbouring terminal nodes, and $\Delta_n=-1$ when a node is exactly at $K$.

\sect{Results used.} Proving the first and the citing the other two results.
\begin{lem}[Prices as tail probabilities]
With $J\sim\mathrm{Bin}(n,p_n)$ and $\hat J\sim\mathrm{Bin}(n,\hat p_n)$, $V_n(g)=c_1S\,\Prob(\hat J\ge a_n)+c_2\,\e^{-rT}\,\Prob(J\ge a_n)$.
\end{lem}
\begin{proof}
By Definitions~2 and~3, $V_n(g)=\e^{-rT}\sum_{j\ge a_n}(c_1Su^jd^{\,n-j}+c_2)\binom nj p_n^jq_n^{\,n-j}$. Since $\e^{-rT}=\e^{-rhj}\e^{-rh(n-j)}$, Definition~1 gives $\e^{-rT}u^jd^{\,n-j}p_n^jq_n^{\,n-j}=\hat p_n^{\,j}(1-\hat p_n)^{n-j}$, and this is the claim used.
\end{proof}
\begin{fact}[Uspensky {[Ch.~VII, p.~129][6]}]
Let $J\sim\mathrm{Bin}(n,p)$, $q=1-p$, $s=\sqrt{npq}\ge5$, and $G(t)=\Phi(t)+\dfrac{q-p}{6s}(1-t^2)\varphi(t)$. Then for all integers $0\le a\le b\le n$,
\[
\Prob(a\le J\le b)=G\Big(\frac{b+\frac12-np}{s}\Big)-G\Big(\frac{a-\frac12-np}{s}\Big)+\Omega, \qquad |\Omega|<\frac{0.13+0.18\,|p-q|}{s^2}+\e^{-3s/2}.
\]
\end{fact}
The shifts by $\tfrac12$ (the continuity correction) reflect that $J$ is integer-valued, and the second term of $G$ its skewness. As it is explicit, the fact applies with $p=p_n$ depending on $n$.
\begin{fact}[Chang-Palmer {[2]}, as restated in {[Thm~4.3 ($\lambda_n=0$), Remark~4.3.1][7]}; and {[3]}]
For the call, $C_n=C+f(\Delta_n)/n+O(n^{-3/2})$, where
\[
f(\Delta)=\frac{S\,\e^{-d_1^2/2}}{24\,\sigma\sqrt{2\pi T}}\Big(A-12\sigma^2T(\Delta^2-1)\Big),\qquad A=-\sigma^2T(6+d_1^2+d_2^2)+4T(d_1^2-d_2^2)\,r-12T^2r^2 .
\]
\end{fact}
\sect{The theorem.}
\begin{thm}
Let $g$ be as in Definition~2. As $n\to\infty$,
\[
V_n(g)=V(g)+\e^{-rT}\varphi(d_2)\,[g]\,\frac{\Delta_n}{\sqrt n}+O\Big(\frac1n\Big).
\]
\textup{(a)} $C_n-C=O(1/n)$ for the call, and \textup{(b)}
$D_n-D=\e^{-rT}\varphi(d_2)\,\Delta_n\,n^{-1/2}+O(1/n)$ for the digital call.
\end{thm}
Part (a) is due to Leisen and Reimer [5]. Part (b) is the leading term of the expansion in [2], and from that I apply Fact~5 to two tail probabilities. As I'm working in the first order, the proof below can be given in full; Fact~6, which goes one order further, is not proved here.
\begin{lem}
With $h=T/n$,
\[
p_n=\frac12+\frac{(r-\sigma^2/2)\sqrt h}{2\sigma}+O(h^{3/2})\qquad\text{and}\qquad \hat p_n=\frac12+\frac{(r+\sigma^2/2)\sqrt h}{2\sigma}+O(h^{3/2}).
\]
\end{lem}
\begin{proof}
Put $s=\sigma\sqrt h$. By Taylor's theorem, $\e^{rh}-\e^{-s}=s+(rh-\tfrac12s^2)+\tfrac16s^3+O(h^2)$ and $u-d=\e^s-\e^{-s}=2s\,(1+\tfrac16s^2+O(h^2))$. Dividing, $p_n=\tfrac12\big(1+\tfrac{rh-s^2/2}{s}+\tfrac16s^2+O(h^{3/2})\big)\big(1-\tfrac16s^2+O(h^2)\big)=\tfrac12+\tfrac{rh-s^2/2}{2s}+O(h^{3/2})$, since the terms $\pm\tfrac16s^2$ cancel. For $\hat p_n$, use $ud=1$ to write $\hat p_n=(u-\e^{-rh})/(u-d)$, where the numerator is $s+(rh+\tfrac12s^2)+\tfrac16s^3+O(h^2)$.
\end{proof}
\begin{lem}[Tail expansion]
Let $\mu\in\R$, let $\pi_n\in(0,1)$ satisfy $\pi_n=\tfrac12+\mu\sqrt h/(2\sigma)+O(h^{3/2})$, and let $J\sim\mathrm{Bin}(n,\pi_n)$. Then, with $\delta=\big(\ln(S/K)+\mu T\big)/(\sigma\sqrt T)$,
\[
\Prob(J\ge a_n)=\Phi(\delta)+\varphi(\delta)\,\frac{\Delta_n}{\sqrt n}+O\Big(\frac1n\Big). 
\]
\end{lem}

\subsect{Central argument: proof of Lemma~9.}
Apply Fact~5 to $[a_n,n]$, remove every term that is $O(1/n)$, and then expand the remaining argument of $\Phi$ to order $n^{-1/2}$. Write $s_n=\sqrt{n\pi_n(1-\pi_n)}$.

\emph{Step 1: Fact~5 applies.} Since $(\pi_n-\tfrac12)^2=O(h)=O(1/n)$, this yields $s_n^2=\tfrac n4-n(\pi_n-\tfrac12)^2=\tfrac n4+O(1)$; so $s_n=\tfrac{\sqrt n}2\,(1+O(1/n))$, and $s_n\ge5$ for large $n$. As $\alpha_n=\tfrac n2+O(\sqrt n)$, also $0\le a_n\le n$ for large $n$. With $t_n=(a_n-\tfrac12-n\pi_n)/s_n$ and $t_n'=(n+\tfrac12-n\pi_n)/s_n$ it follows that $\Prob(J\ge a_n)=G(t_n')-G(t_n)+\Omega_n$ with $\Omega_n=O(1/n)$.

\emph{Step 2: the terms that are $O(1/n)$.} Since $|(1-t^2)\varphi(t)|\le1$ and $|1-2\pi_n|=O(n^{-1/2})$, this gives $|G(t)-\Phi(t)|\le|1-2\pi_n|/(6s_n)=O(1/n)$ for every $t$. Here I need $\pi_n\to\tfrac12$ at rate $n^{-1/2}$; for a fixed $\pi\ne\tfrac12$ the skewness term would be of order $n^{-1/2}$.
Next, $t_n'\ge n(1-\pi_n)/s_n=\sqrt{n(1-\pi_n)/\pi_n}\ge\sqrt{n/2}$ once $\pi_n\le\tfrac23$, and for $t>0$, $1-\Phi(t)=\int_t^\infty\varphi(x)\,dx\le\int_t^\infty\tfrac xt\varphi(x)\,dx=\varphi(t)/t$, so $1-\Phi(t_n')=O(\e^{-n/4})$. Thus, $\Prob(J\ge a_n)=1-\Phi(t_n)+O(1/n)=\Phi(-t_n)+O(1/n)$.

\emph{Step 3: expansion of $t_n$; where $\Delta_n$ appears.} Write $a_n=\alpha_n+\theta_n$ and $\theta_n-\tfrac12=-\tfrac12\Delta_n$. Then $a_n-\tfrac12-n\pi_n=\ln(K/S)/(2\sigma\sqrt h)-n(\pi_n-\tfrac12)+(\theta_n-\tfrac12)$, and since $1/\sqrt h=\sqrt{n/T}$, $n\sqrt h=\sqrt{nT}$ and $n\,O(h^{3/2})=O(n^{-1/2})$,
\[
a_n-\tfrac12-n\pi_n =\frac{\sqrt n}{2}\cdot\frac{\ln(K/S)-\mu T}{\sigma\sqrt T}-\frac{\Delta_n}{2}+O(n^{-1/2}) =-\frac{\sqrt n}{2}\,\delta-\frac{\Delta_n}{2}+O(n^{-1/2}).
\]
Dividing by $s_n=\tfrac{\sqrt n}{2}(1+O(1/n))$ gives $t_n=-\delta-\Delta_n/\sqrt n+O(1/n)$. This step uses two properties of $\pi_n$: its first-order coefficient, which gives $\delta$, and the size $O(h^{3/2})$ of its remainder. An error of order $h$ in $\pi_n$ would add $n\cdot O(h)=O(1)$ here, i.e.\ $O(n^{-1/2})$ in $t_n$, and break the result. The continuity correction $\tfrac12$ turns $\theta_n$ into $-\tfrac12\Delta_n$.

\emph{Step 4: Taylor expansion.} $\Phi'=\varphi$ and $\Phi''(t)=-t\varphi(t)$ are bounded and $|\Delta_n|\le1$, so $\Phi(-t_n)=\Phi\big(\delta+\Delta_n/\sqrt n+O(1/n)\big)=\Phi(\delta)+\varphi(\delta)\Delta_n/\sqrt n+O(1/n)$. With Step 2 this proves the lemma.\qed

\begin{proof}[Proof of Theorem~7]
By Lemma~8, Lemma~9 applies to $\pi_n=p_n$ with $\mu=r-\sigma^2/2$, $\delta=d_2$, and to $\pi_n=\hat p_n$ with $\mu=r+\sigma^2/2$, $\delta=d_1$. Since both tails in Lemma~4 start at the same node $a_n$, they carry the same $\Delta_n$:
\[
V_n(g)=c_1S\,\Phi(d_1)+c_2\e^{-rT}\Phi(d_2)+\big(c_1S\varphi(d_1)+c_2\e^{-rT}\varphi(d_2)\big)\frac{\Delta_n}{\sqrt n}+O\Big(\frac1n\Big).
\]
Now $d_1^2-d_2^2=(d_1-d_2)(d_1+d_2)=2\ln(S/K)+2rT$, so $\varphi(d_1)/\varphi(d_2)=\e^{-(d_1^2-d_2^2)/2}=K\e^{-rT}/S$, that is
$S\varphi(d_1)=K\e^{-rT}\varphi(d_2)$ (the identity used for the Greeks in [proof of Prop.~7.9][1]). Hence the bracket
equals $\e^{-rT}\varphi(d_2)(c_1K+c_2)=\e^{-rT}\varphi(d_2)[g]$. For the call, $[g]=0$, so the $n^{-1/2}$ terms cancel.
\end{proof}







\begin {thebibliography}{9}
\bibitem{bandtlow}O.~Bandtlow, \emph{MTH6154 Financial Mathematics I}, lecture notes, Queen Mary University of London, Semester A 2024.
\bibitem{changpalmer}L.-B.~Chang, K.~Palmer, Smooth convergence in the binomial model, \emph{Finance Stoch.}\ 11 (2007), 91--105.
\bibitem{diener}F.~Diener, M.~Diener, Asymptotics of the price oscillations of a European call option in a tree model, \emph{Math.\ Finance} 14 (2004), 271--293.
\bibitem{feller}W.~Feller, \emph{An Introduction to Probability Theory and Its Applications}, Vol.~II, 2nd ed., Wiley, 1971.
\bibitem{leisen}D.~Leisen, M.~Reimer, Binomial models for option valuation: examining and improving convergence, \emph{Appl.\ Math.\ Finance} 3 (1996), 319--346.
\bibitem{uspensky}J.~V.~Uspensky, \emph{Introduction to Mathematical Probability}, McGraw--Hill, 1937.
\bibitem{viitasaari}L.~Viitasaari, Rate of convergence for discrete approximation of option prices, arXiv:1207.6756v4, 2013.
\end{thebibliography}



\end{document}